\documentclass[12pt,a4paper]{article}

\usepackage[utf8]{inputenc}
\usepackage[T1]{fontenc}
\usepackage[spanish]{babel}
\usepackage{amsmath, amssymb}
\usepackage{geometry}
\usepackage{setspace}
\usepackage{parskip}
\usepackage{enumitem}
\usepackage{hyperref}

\geometry{margin=2.5cm}
\onehalfspacing

\begin{document}

\begin{center}
{\Large\bfseries Propuesta de hitos para Proyecto de T\'itulo II}
\end{center}

\vspace{0.3cm}
\noindent\textbf{Nombre Alumno:} Eduardo Caballero\\
\textbf{Nombre Profesor gu\'ia:} Sebasti\'an Cea\\
\textbf{Tema Memoria:} Modelaci\'on de redes de cr\'edito bajo informaci\'on asim\'etrica.

\section*{HITO I}

\subsection*{Objetivo:}

El objetivo de este hito es formalizar y resolver anal\'iticamente el modelo heterog\'eneo
propuesto en el Hito~2, relajando los dos supuestos de homogeneidad del marco de Venegas
(2024) ($\gamma_i = 1$ para todo agente y $\mu_i = \mu_j$) y caracterizando el equilibrio
\emph{pairwise-stable} del sistema de negociaci\'on bilateral bajo un coeficiente de
aversi\'on al riesgo heterog\'eneo $\Gamma$. Se busca demostrar consistencia del modelo
extendido con los resultados de Venegas (2024) en los casos l\'imite correspondientes, e
implementar y validar num\'ericamente el modelo extendido contra los benchmarks ya
publicados, sentando la base cuantitativa sobre la cual se apoyan las simulaciones del
Hito~II.

\subsection*{Evaluaci\'on:}

El cumplimiento de este hito ser\'a evaluado a trav\'es de la entrega de un documento que
incluya:
\begin{itemize}
    \item Formalizaci\'on matem\'atica del modelo con agentes heterog\'eneos en aversi\'on al
    riesgo ($\gamma_i$) y creencias sobre retornos ($\mu_i$).
    \item Caracterizaci\'on anal\'itica del equilibrio \emph{pairwise-stable} bajo $\Gamma$
    heterog\'eneo, con verificaci\'on de consistencia en los casos l\'imite del modelo de
    Venegas (2024) ($\lambda_A \to 0$, $\tau = 0$, $\Psi^A$ inactiva, $\Psi^T = \Psi^A$).
    \item Implementaci\'on computacional del modelo extendido y validaci\'on num\'erica contra
    los benchmarks ya publicados de Venegas (2024) (3F-WITH/WITHOUT, 4F-WITH/WITHOUT).
\end{itemize}

\section*{HITO II}

\subsection*{Objetivo:}

El objetivo de este hito es responder, mediante simulaci\'on, las tres preguntas de
investigaci\'on formuladas en el Hito~2 sobre el modelo heterog\'eneo validado en el Hito~I, y
dar inicio a la aplicaci\'on emp\'irica del modelo. Esto incluye evaluar la sensibilidad de los
resultados de bienestar y comportamiento estrat\'egico ante variaciones en los par\'ametros del
modelo, as\'i como seleccionar y preparar un caso de datos reales sobre el cual calibrar
preliminarmente el modelo extendido, siguiendo el precedente metodol\'ogico de Jalan \&
Chakrabarti (2024) en su validaci\'on emp\'irica con datos de comercio internacional.

\subsection*{Evaluaci\'on:}

El cumplimiento de este hito ser\'a evaluado a trav\'es de la entrega de un documento que
incluya:
\begin{itemize}
    \item Resultados de simulaci\'on para las tres preguntas de investigaci\'on: persistencia de
    la recuperaci\'on de bienestar, \emph{sorting} end\'ogeno entre redes, y distribuci\'on de
    las ganancias de bienestar entre agentes.
    \item An\'alisis de sensibilidad de los resultados ante variaciones en
    $(\gamma_i, \mu_i, \tau, \lambda_A)$.
    \item Selecci\'on del caso de datos reales para la aplicaci\'on emp\'irica, con el dataset
    identificado y una primera calibraci\'on preliminar del modelo sobre ese caso.
\end{itemize}

\end{document}
